% TOP_PROBLEM: {"id": 3700026, "title": "Real two-one-point extremals", "category_id": 8, "category": {"id": 8, "name": "analysis", "display_name": "Analysis", "description": "Limits, continuity, calculus, and function theory.", "slug": "analysis", "order_index": 8, "created_at": "2026-07-31T15:26:25.671Z"}, "status": "partially_solved", "problem_number": "AMR-036-0026", "set_id": 13, "statusQualification": "The source's 'same questions' retains the simple zero and multiplicity-counting conventions from Problem 1."}
% FIRST_POSTED: 2026-10-01
% ENTRY_KIND: statement_only
% UnsolvedMath ID 3700026; source catalogue code AMR-036-0026
% !TeX program = lualatex
% Definition and sources only; this file is not an LLM solution attempt.
\documentclass[11pt,a4paper]{article}
\usepackage[margin=25mm]{geometry}
\usepackage{fontspec}
\setmainfont{lmroman10-regular.otf}[BoldFont=lmroman10-bold.otf,
 ItalicFont=lmroman10-italic.otf,BoldItalicFont=lmroman10-bolditalic.otf]
\usepackage{amsmath,amssymb,mathtools}
\usepackage{unicode-math}
\setmathfont{latinmodern-math.otf}
\AtBeginDocument{\let\setminus\smallsetminus}
\usepackage{xurl,hyperref}
\hypersetup{colorlinks=true,urlcolor=blue}
\setcounter{secnumdepth}{0}
\setlength{\parindent}{0pt}
\setlength{\parskip}{5pt}
\setlength{\emergencystretch}{3em}
\begin{document}
\section*{Real two-one-point extremals}\label{3700026}
{\small UnsolvedMath ID 3700026 · AMR-036-0026 · Analysis · AMR Open Problem Lists}
\subsection{Definitions and mathematical statement}
Call a holomorphic function on $\mathbb D$ real if
$f(\overline z)=\overline{f(z)}$ for every $z\in\mathbb D$,
equivalently if its Taylor coefficients at the origin
are real. Let $\mathcal F_{\mathbb R}$ consist of these
functions with one simple zero at $0$, no other zeros
in the disk, and exactly two $1$-points $z_1,z_2$,
counted with multiplicity.

Determine
\[
 R_{\mathbb R}=\inf_{f\in\mathcal F_{\mathbb R}}
                         \max\{|z_1|,|z_2|\},
 \qquad
 D_{\mathbb R}=\sup_{f\in\mathcal F_{\mathbb R}}|f'(0)|.
\]
For each optimization, decide whether its sharp value
is attained and describe every extremal function, or
appropriate extremizing sequences if necessary.
The derivative is a nonzero real number, but its
absolute value is the objective.

Reality forces the $1$-point multiset to be invariant
under conjugation. It may therefore consist of two
real points, a nonreal conjugate pair, or a single
real point of multiplicity two. All these possibilities
are included; requiring simple or nonreal $1$-points
would give a different extremal problem.
No univalence or bounded-range hypothesis is made,
and the functions need not extend across the unit
circle. The zero remains fixed at the center and
the radius is Euclidean. The aim is the exact effect
of the reality constraint on both sharp quantities
and their equality cases.
\subsection{Short English statement}
Determine the two-one-point radius and derivative extremals when the disk function has real Taylor coefficients.
\subsection{Sources}
\begin{itemize}
\item[S1] ULAM.AI and the UnsolvedMath Contributors, supplied problems.json, record 3700026 (AMR-036-0026), AMR Open Problem Lists. Catalogue identity and source transcription; CC BY 4.0.
\url{https://huggingface.co/datasets/ulamai/UnsolvedMath}.
\item[S2] Eremenko - My favorite unsolved problems, Goldberg relatives, Problem 2. Original source indexed by AMR-036-0026.
\url{https://www.math.purdue.edu/~eremenko/uns1.html}.
\item[S3] A. Eremenko, Goldberg's constant and its relatives, Problem 2. Author-maintained primary problem note.
\url{https://www.math.purdue.edu/~eremenko/dvi/goldconst.pdf}.
\end{itemize}
\end{document}
